\problemname{Skogsvaktare}
Du har ett träd, d.v.s. en sammanhängande graf utan cykler.
På vissa av dessa hörn växer det frukter som du vill vakta mot inkräktande maskar och annan ohyra.
Till din hjälp har du designat en insektsrobot som kan vakta frukter.
Om du placerar en insektsrobot på ett visst hörn så vaktar den alla de frukter vars avstånd från roboten är så litet som möjligt.
Mer formellt, antag att det minsta avståndet från roboten till någon frukt är $D$.
Då vaktar roboten alla frukter som har exakt avståndet $D$ från roboten.

Du vill placera robotar i trädet så att du vaktar samtliga frukter, men av kostnadsskäl vill du minimera antalet robotar du placerar.
Finn en sådan placering av robotar i trädet.

\section*{Indata}
Den första raden av indatan består av två heltal $N$ och $F$ ($1 \le N \le 500\,000$, $1 \le F \le N$), antalet hörn och antalet frukter i trädet.

De följande $N - 1$ raderna beskriver kanterna i trädet.
Hörnen i trädet är numrerade mellan $1$ och $N$.
En kant ges av två heltal $a$ och $b$ ($1 \le a, b \le N$) och betyder att kanten förbinder hörn $a$ och $b$.

Till slut följer en rad med $F$ olika heltal, de hörn där det finns frukter.

\section*{Utdata}
Skriv först ut det minsta antalet hörn $K$ du måste placera robotar på.

Skriv sedan ut en rad med $K$ heltal på, hörnen som du placerar robotar på.

\section*{Poängsättning}
Din lösning kommer att testas på en mängd testfallsgrupper.
För att få poäng för en grupp så måste du klara alla testfall i gruppen.

\noindent
\begin{tabular}{| l | l | p{12cm} |}
  \hline
  \textbf{Grupp} & \textbf{Poäng} & \textbf{Gränser} \\ \hline
    $1$   & $8$      & Den $i$:te kanten förbinder hörnen $i$ och $i + 1$, d.v.s. trädet är en linje. \\ \hline
    $2$   & $18$     & $F \le 15$ \\ \hline
    $3$   & $23$     & $N \le 2\,000$ \\ \hline
    $4$   & $51$     & Inga ytterligare begränsningar. \\ \hline
\end{tabular}

